\chapter*{چطور این کتابچه را بخوانیم}
\addcontentsline{toc}{chapter}{چطور این کتابچه را بخوانیم}
\markboth{چطور این کتابچه را بخوانیم}{چطور این کتابچه را بخوانیم}

این کتابچه بخشی مستقل و خودکفاست --- یک تکه‌ی کامل و قابل‌خواندن از
یک اثر بزرگ‌تر، نه پاره‌ای که برای معنا پیدا کردن به بقیه‌ی کتاب نیاز
داشته باشد. هر چیزی که برای دنبال کردن استدلال از اولین صفحه تا آخرین
صفحه لازم است، همین‌جاست.

چند قرارداد از پروژه‌ی اصلی به اینجا هم رسیده‌اند و پیش از فصل اول
ارزش دانستن دارند:

\begin{notebox}[title=یادداشت]
بخش‌هایی که به این شکل جدا شده‌اند، چیزی را نشان می‌دهند که متن اصلی
بدون توضیح بیشتر روی آن تکیه می‌کند: اصطلاحی که ارزش دقیق‌تر شدن دارد،
تمایزی که به‌راحتی محو می‌شود، یا نکته‌ای حاشیه‌ای که اهمیتش بیشتر از
حجمش است.
\end{notebox}

\begin{tipbox}[title=نکته]
این سبک یک پیشنهاد عملی را نشان می‌دهد --- روشی متفاوت برای خواندن یک
بخش، یا اشاره‌ای به چیزی که بعد از تمام شدن فصل ارزش دنبال کردن دارد.
\end{tipbox}

پانویس‌ها و ارجاع‌های داخل متن به فصل منابع در انتهای کتابچه اشاره
می‌کنند، جایی که همه‌ی آثاری که واقعاً در این کتابچه به آن‌ها استناد
شده، به‌طور کامل فهرست شده‌اند. بخش‌هایی که به‌عنوان «در دست نوشتن»
علامت‌گذاری شده‌اند دقیقاً همین‌اند: استدلال تا همان نقطه روی پای خودش
می‌ایستد، اما ممکن است خودِ آن بخش در ویرایش بعدی کامل‌تر شود.

فراتر از این، ترتیبی جز همان ترتیب فصل‌ها الزامی نیست، و پیش‌زمینه‌ای
فرض نشده که این کتابچه خودش در طول مسیر آن را نساخته باشد.
